\documentclass[../main.tex]{subfiles}
\begin{document}

\section{Đặt vấn đề}
\label{sec:dvd}

Các hệ thống điều khiển công nghiệp giữ vai trò quan trọng trong nhiều lĩnh vực của đời sống như sản xuất, năng lượng, xử lý nước, giao thông và hạ tầng kỹ thuật. Do đặc thù vận hành liên tục, lâu dài và yêu cầu cao về độ ổn định, các hệ thống công nghiệp truyền thống thường được thiết kế với trọng tâm là tính sẵn sàng, độ tin cậy và khả năng đáp ứng thời gian thực. Mặt trái của việc để đạt được các ưu tiên trên, các giao thức công nghiệp được thiết kế mang tính độc quyền/mở nhưng bị kiểm soát chặt chẽ để gắn chặt với phần cứng của các tập đoàn công nghiệp nặng lâu đời để đảm bảo tính tin cậy, ví dụ như S7 Comm của Siemens (Đức), MC của Mitsubishi Electric (Nhật), không có các cơ chế mã hoá, xác thực hay toàn vẹn để đảm bảo tính thời gian thực. Yêu cầu về bảo mật đặt ra cho các giao thức này ban đầu là không cần thiết do mạng OT của nhà máy thường hoạt động khép kín với thế giới bên ngoài nên bề mặt tấn công là gần như không có với kẻ tấn công từ bên ngoài. Tuy nhiên, để quản lý quy mô sản xuất ngày càng phức tạp của các nhà máy, yêu cầu bắt buộc đặt ra là mạng OT phải được kết nối đến các hệ thống giám sát tập trung (SCADA), tới hệ thống quản lý sản xuất (MES) và tới mạng IT của doanh nghiệp để điều phối hoạt động sản xuất tương thích với cá hoạt động khác như tài chính, chuỗi cung ứng,.. (ERP).Điều này đã làm thay đổi bề mặt tấn công của môi trường OT. Khi này các giao thức được sử dụng trong môi trường có mức độ kết nối cao hơn, kẻ tấn công có thêm bề mặt để truy cập được tới mạng OT của nhà máy.

Hiện tại nhà máy tại Việt Nam và các dự án trọng điểm quốc gia đều triển khai sử dụng thiết bị do tập đoàn Siemens của Đức chế tạo. Giao tiếp giữa các thiết bị của Siemens gần như được thực hiện thông qua giao thức độc quyền S7Comm hoặc bản cải tiến S7Comm Plus. Giao thức này cho phép các thành phần như Simantic HMI, phần mềm lập trình Tia Portal, PLC S7-300/400/1200/1500, biến tần và các thiết bị khác trong hệ sinh thái Siemens trao đổi dữ liệu phục vụ giám sát, cấu hình và điều khiển. Tuy nhiên, trong các cấu hình S7 Comm truyền thống, các thao tác như đọc dữ liệu, ghi dữ liệu hoặc gửi lệnh điều khiển có thể được thực hiện chỉ cần thiết bị trong mạng có đủ thông tin cần thiết về địa chỉ và cấu trúc bộ nhớ của PLC. Trong môi trường nhà máy, hai thông tin trên không phải là ưu tiên bảo mật và vì thế có thể dễ dàng được truy cập bởi bất cứ kỹ sư nào trong xưởng. Điều này tạo ra nguy cơ xuất hiện các hành vi tấn công như dò quét thiết bị, thu thập thông tin hệ thống, ghi trái phép giá trị quá trình, làm quá tải dịch vụ S7 (DoS), phát lại lệnh điều khiển hoặc can thiệp vào luồng truyền thông giữa PLC và trạm giám sát được thực hiện mà không gặp bất cứ khó khăn đáng kể nào.

Bài toán đặt ra là cần có một phương pháp hỗ trợ phát hiện các hành vi bất thường trong truyền thông công nghiệp mà không làm ảnh hưởng đến hoạt động của hệ thống điều khiển do đây là yêu cầu cao nhất của nhà máy. Mọi hành động đảm bảo an ninh nhưng gây chậm dây chuyền hoặc nguy hiểm hơn là dừng dây chuyền sẽ đều là không có giá trị ứng dụng thực tiễn. Do đó, việc nghiên cứu một hướng tiếp cận giám sát thụ động dựa trên lưu lượng mạng có ý nghĩa thực tiễn trong việc nâng cao khả năng quan sát an toàn thông tin cho hệ thống công nghiệp mà không gây làm gián đoạn đến hoạt động sản xuất.

\section{Tổng quan và các vấn đề nghiên cứu}
\label{sec:giaiphap}

Các giải pháp bảo vệ hệ thống điều khiển công nghiệp hiện nay thường được triển khai theo nhiều lớp. Ở lớp kiến trúc mạng, hệ thống có thể được phân vùng theo mô hình Purdue, sử dụng tường lửa công nghiệp, vùng DMZ và các cơ chế kiểm soát truy cập để hạn chế kết nối không cần thiết giữa mạng IT và mạng OT. Ở lớp thiết bị hiện trường, controller đặt tại nhà máy, các biện pháp như cấu hình quyền truy cập, cập nhật phần mềm, tắt dịch vụ không sử dụng và giới hạn chức năng truyền thông có thể giúp giảm nguy cơ bị khai thác. Ở lớp giám sát, các hệ thống giám sát tập trung như SCADA được sử dụng để giám sát các thông số vận hành hay hệ thống phát hiện xâm nhập IDS sử dụng để phân tích lưu lượng mạng, từ đó nhận diện mẫu tấn công và cảnh báo khi xuất hiện hành vi bất thường.

Trong đó, hệ thống phát hiện xâm nhập dựa trên mạng là một hướng tiếp cận phù hợp với môi trường công nghiệp. Thay vì can thiệp trực tiếp vào các thiết bị đang vận hành trong dây chuyền, hệ thống phát hiện xâm nhập có thể được triển khai ở chế độ thụ động thông qua cơ chế sao chép lưu lượng từ switch sử dụng SPAN port/Mirror Port hoặc các thiết bị phần cứng chuyen dụng như TAP. Cách triển khai này giúp hệ thống giám sát quan sát được các trao đổi giữa các thành phần trong mạng mà không làm thay đổi thay đổi lưu lượng mạng, cấu hình các thiết bị hiện có cũng như việc có thể triển khai trực tiếp trong quá trình dây chuyền đang vận hành mà không gây Downtime.

Một số công cụ mã nguồn mở như Suricata có thể được sử dụng để xây dựng cơ chế phát hiện dựa trên chữ ký. Với cách tiếp cận này, các luật phát hiện được xây dựng dựa trên đặc điểm của giao thức, mã chức năng, cấu trúc gói tin hoặc chuỗi byte đặc trưng trong lưu lượng mạng. Đối với giao thức S7comm, các trường như TPKT, COTP, mã nhận dạng giao thức S7 và mã chức năng có thể được sử dụng để nhận diện các thao tác đọc, ghi, truyền chương trình hoặc điều khiển trạng thái PLC. Phương pháp này có ưu điểm là rõ ràng, dễ kiểm chứng và phù hợp với các kịch bản tấn công đã biết.

Tuy nhiên, các giải pháp phát hiện dựa trên chữ ký cũng tồn tại một số hạn chế. Thứ nhất, hệ thống chỉ có thể phát hiện tốt các hành vi đã được mô tả trong luật, trong khi các biến thể tấn công mới hoặc các hành vi bất thường ở mức logic quá trình có thể khó nhận diện nếu chỉ dựa vào mẫu gói tin cố định. Thứ hai, trong môi trường công nghiệp, một số thao tác hợp lệ về mặt kỹ thuật vẫn có thể gây rủi ro nếu xuất hiện sai thời điểm, sai nguồn gửi hoặc với tần suất bất thường. Thứ ba, việc xây dựng luật phát hiện cần hiểu rõ cấu trúc giao thức và đặc điểm vận hành của hệ thống cụ thể để hạn chế cảnh báo sai.

Từ các vấn đề trên, đồ án tập trung nghiên cứu hướng phát hiện các hành vi tấn công điển hình trên giao thức S7comm trong môi trường mô phỏng. Phạm vi nghiên cứu không nhằm xây dựng một hệ thống bảo vệ hoàn chỉnh cho mọi loại mạng công nghiệp, mà tập trung vào việc mô phỏng, phân tích và xây dựng các luật phát hiện cho một số hành vi có ý nghĩa đại diện. Cách tiếp cận này giúp làm rõ mối liên hệ giữa điểm yếu của giao thức, biểu hiện của tấn công trong lưu lượng mạng và khả năng phát hiện bằng hệ thống IDS thụ động.

\section{Mục tiêu nghiên cứu và định hướng giải pháp}

Mục tiêu chính của đồ án là nghiên cứu và xây dựng một giải pháp phát hiện xâm nhập thụ động cho một số hành vi tấn công trên giao thức S7comm trong môi trường điều khiển công nghiệp mô phỏng. Để đạt được mục tiêu này, đồ án trước hết cần tìm hiểu về khái niệm của OT, đặc điểm của hệ thống điều khiển công nghiệp, cấu trúc mạng Purdue, vai trò của PLC và cách thức giao tiếp của giao thức S7comm. Trên cơ sở đó, đồ án phân tích các nguy cơ bảo mật có thể phát sinh khi giao thức công nghiệp thiếu cơ chế xác thực và bảo vệ dữ liệu đầy đủ.

Định hướng giải pháp của đồ án là xây dựng môi trường mô phỏng có khả năng tạo ra cả lưu lượng hợp lệ và lưu lượng tấn công. Các kịch bản tấn công được lựa chọn nhằm phản ánh một số nhóm hành vi phổ biến trong quá trình khai thác hệ thống PLC, bao gồm thu thập thông tin, thao tác dữ liệu quá trình, từ chối dịch vụ trên kênh S7comm, phát lại lệnh điều khiển và can thiệp vào truyền thông. Việc mô phỏng các kịch bản này giúp tạo cơ sở thực nghiệm để quan sát cấu trúc gói tin, xác định dấu hiệu nhận biết và đánh giá khả năng phát hiện.

Sau khi có dữ liệu lưu lượng, đồ án sử dụng hệ thống IDS Suricata để xây dựng các luật phát hiện tương ứng: \texttt{s7comm.rules} cho từng opcode/payload đặc trưng và \texttt{ics\_dos.rules} cho flood (ngưỡng tần suất). Giải pháp được triển khai theo hướng thụ động, trong đó hệ thống phát hiện chỉ quan sát lưu lượng mạng và sinh cảnh báo khi phát hiện mẫu hành vi phù hợp. Cách tiếp cận này phù hợp với yêu cầu của môi trường công nghiệp do các thiết bị OT thường không nhận được các bản cập nhật, được thi công bởi nhiều nhà thầu từ nhiều quốc gia, việc thay đổi cấu hình hoạt động mang lại rủi ro lớn trong việc tương thích cũng như khả năng khôi phục vận hành.

Bên cạnh việc xây dựng luật phát hiện, đồ án cũng hướng tới việc đánh giá khả năng hoạt động của các luật trong từng kịch bản mô phỏng. Quá trình đánh giá tập trung vào việc kiểm tra hệ thống có nhận diện đúng các hành vi tấn công đã mô phỏng hay không, đồng thời xem xét các tình huống có thể gây cảnh báo sai. Từ đó, đồ án rút ra nhận xét về ưu điểm, hạn chế và phạm vi áp dụng của phương pháp phát hiện dựa trên chữ ký đối với giao thức S7comm.

\section{Đóng góp của đồ án}

Đồ án này có các đóng góp chính như sau:

\begin{enumerate}
    \item Đồ án trình bày tổng quan về lĩnh vực OT/ICS, một số vấn đề an toàn thông tin trong hệ thống điều khiển công nghiệp, tập trung vào rủi ro phát sinh từ giao thức S7comm trong môi trường truyền thông PLC.

    \item Đồ án xây dựng môi trường mô phỏng phục vụ nghiên cứu các hành vi tấn công điển hình trên giao thức S7comm, bao gồm trinh sát, thao tác dữ liệu (ghi DB), DoS S7comm (script \texttt{s7comm\_dos.py}), phát lại Start/Stop và quan sát MITM.

    \item Đồ án phân tích đặc điểm lưu lượng mạng của các kịch bản tấn công đã mô phỏng, từ đó xác định các dấu hiệu có thể sử dụng để phát hiện bằng hệ thống IDS.

    \item Đồ án xây dựng tập luật phát hiện dựa trên Suricata cho một số hành vi tấn công S7comm, đồng thời kiểm tra khả năng sinh cảnh báo của các luật trong môi trường thực nghiệm.

    \item Đồ án đưa ra nhận xét về khả năng áp dụng, ưu điểm và hạn chế của phương pháp phát hiện xâm nhập thụ động dựa trên chữ ký trong bối cảnh mạng điều khiển công nghiệp.
\end{enumerate}

\section{Bố cục đồ án}

Phần còn lại của báo cáo đồ án tốt nghiệp được tổ chức như sau.

Chương 2 trình bày cơ sở lý thuyết liên quan đến hệ thống điều khiển công nghiệp, kiến trúc mạng OT, PLC và giao thức truyền thông S7comm. Chương này cũng giới thiệu các khái niệm nền tảng về an toàn thông tin trong môi trường công nghiệp và hệ thống phát hiện xâm nhập dựa trên mạng. Nội dung của chương là cơ sở để phân tích các nguy cơ bảo mật và xây dựng giải pháp phát hiện trong các chương tiếp theo.

Chương 3 trình bày quá trình xây dựng môi trường mô phỏng và các kịch bản tấn công được sử dụng trong đồ án. Các kịch bản được lựa chọn nhằm thể hiện một số hành vi có rủi ro đối với hệ thống PLC, bao gồm thu thập thông tin, thay đổi dữ liệu quá trình, DoS trên kênh S7comm (Mục~\ref{sec:dos_s7comm}), phát lại lệnh điều khiển và can thiệp vào luồng truyền thông. Thông qua chương này, đồ án làm rõ cách tạo dữ liệu thực nghiệm phục vụ cho việc phân tích và phát hiện.

Chương 4 trình bày phương pháp phát hiện các hành vi tấn công S7comm bằng hệ thống IDS Suricata: phân tích cấu trúc gói tin, lựa chọn dấu hiệu và xây dựng luật trong \texttt{s7comm.rules} cùng \texttt{ics\_dos.rules} (Mục~\ref{sec:rules-dos}).

Chương 5 trình bày thiết lập đánh giá lab và kết quả kiểm thử, trong đó có phát hiện DoS qua \texttt{sid:1000040}--\texttt{1000043} (Mục~\ref{sec:eval-dos}). Chương 6 tổng kết, nêu hạn chế và hướng phát triển (hiệu chỉnh ngưỡng DoS, giảm cảnh báo sai, mở rộng kịch bản).

\end{document}